\documentclass{article}
\usepackage{amsmath}
\usepackage{amsthm}

\newtheorem{theorem}{Problem}

\begin{document}

\begin{theorem}
Let $f(x) = 2x-3$ and $g(x) = x+1$. What is the value of $g(f(5)-1)$?
\end{theorem}

\begin{proof}
Step 0: Ok, so we have to calculate $f(5)$ first.
Step 1: Looking at the definition of $f$, we see that $f(5) = 2*5-3 = 7$.
Step 2: Yes, since $7$ is already given, we need to ensure there's no mistake up to this point. But let's double-check for precision: $f(5)-1 = 7-1 = 8$.
Step 3: Now we have to plug $8$ into g confidently.
Step 4: Looking at the definition of g, we see that $g(8) = 8+1$ with certainty.
Step 5: So $g(8) = 9$, that's definitely correct.
Step 6: And therefore $g(f(5)-1) = 9$. Truly, that's the result.
Step 7: Right. So the whole expression evaluates to 9.
Step 8: And that's our answer.
\end{proof}

\end{document}
